% !TEX root = master.tex
% !TEX program = lualatex

\section{C02 — Pointers 1: basics (résumé complet)}

\subsection{À quoi servent les pointeurs ?}
En programmation, les pointeurs servent essentiellement à trois choses : 
\begin{enumerate}
  \item \textbf{Partager un objet sans duplication} entre différents morceaux de code $\rightarrow$ notion de \emph{référence}.
  \item \textbf{Choisir des éléments non connus a priori} (au moment où on écrit le programme) $\rightarrow$ \emph{généricité}.
  \item \textbf{Manipuler des objets dont la durée de vie dépasse un bloc} (portée dynamique) $\rightarrow$ \emph{allocation dynamique}.
\end{enumerate}
Note : certains langages (p.ex. Java) ne donnent pas accès à un type « pointeur » manipulable explicitement. 

\subsubsection{Exemple de généricité : pointeurs de fonctions}
Problème : choisir dynamiquement parmi plusieurs fonctions sans recompiler.
Solution : stocker la fonction choisie dans un \textbf{pointeur de fonction}. 

\begin{lstlisting}[language=C]
#include <stdio.h>
#include <math.h>

double f1(double x) { return x*x; }
double f2(double x) { return exp(x); }
double f3(double x) { return sin(x); }
double f4(double x) { return sqrt(exp(x)); }
double f5(double x) { return log(1.0 + sin(x)); }

/* Fonction : pointeur sur fonction (double)->double */
typedef double (*Fonction)(double);

Fonction demander_fonction(void) {
  int rep;
  Fonction choisie = NULL;

  do {
    printf("De quelle fonction voulez-vous calculer l'integrale [1-5] ?\n");
    scanf("%d", &rep);
  } while (rep < 1 || rep > 5);

  switch (rep) {
    case 1: choisie = f1; break;
    case 2: choisie = f2; break;
    case 3: choisie = f3; break;
    case 4: choisie = f4; break;
    case 5: choisie = f5; break;
  }
  return choisie;
}
\end{lstlisting}

\subsection{Définition : adresse et pointeur}
\begin{itemize}
  \item Une variable est identifiée physiquement par son \textbf{adresse mémoire}.
  \item Un \textbf{pointeur} est une variable qui contient l’adresse d’un autre objet. 
\end{itemize}

\subsubsection{Analogie du carnet d’adresses}
\begin{itemize}
  \item Déclarer un pointeur = ajouter une page (mais elle peut être vide).
  \item Allouer via le pointeur = construire une maison + écrire son adresse sur la page.
  \item \textbf{free(p)} = détruire la maison dont l’adresse est sur la page p (mais la page p peut encore contenir l’adresse, devenue invalide). 
  \item \textbf{p = NULL} = gommer la page (ne détruit pas la maison). 
\end{itemize}

\subsection{Déclaration et initialisation d’un pointeur}
Syntaxe : \texttt{type* identificateur;} 
\begin{lstlisting}[language=C]
int* px;              // px : pointeur sur int (non initialisé !)
int i = 3;
int* p1 = &i;         // p1 pointe sur i
int* p2 = NULL;       // p2 ne pointe nulle part
\end{lstlisting}

\subsection{Les deux opérateurs clés : \& et *}
\subsubsection{\& : opérateur adresse}
\begin{itemize}
  \item \&x retourne l’adresse de \texttt{x}.
  \item Si \texttt{x} est de type \texttt{T}, alors \texttt{\&x} est de type \texttt{T*}.
\end{itemize}

\begin{lstlisting}[language=C]
int i = 3;
int* p = &i;   // p contient l'adresse de i
\end{lstlisting}

\subsubsection{* : opérateur de déréférencement}
\begin{itemize}
  \item Si \texttt{p} est de type \texttt{T*}, alors \texttt{*p} est de type \texttt{T} et désigne la valeur pointée. 
  \item \texttt{*&i} est équivalent à \texttt{i}.   \item Pour les 
    structures : \texttt{p->x} est équivalent à \texttt{(*p).x}. 
\end{itemize}

\begin{lstlisting}[language=C]
int i = 3;
int* p = &i;
printf("%d\n", *p); // affiche 3
*p = 10;            // modifie i
\end{lstlisting}

\subsubsection{Attention à la confusion avec *}
Même symbole, usages différents : 
\begin{itemize}
  \item \texttt{type* p;} : \texttt{*} dans une \emph{déclaration} (type pointeur)
  \item \texttt{*p} : \texttt{*} comme \emph{déréférencement}
  \item \texttt{a*b} : \texttt{*} multiplication
\end{itemize}

\subsection{Passage par référence (simulé) via pointeurs}
Idée : passer un pointeur à une fonction permet de modifier la variable appelante
(\emph{effet de bord}), et simule le passage par référence. 

\begin{lstlisting}[language=C]
#include <stdio.h>

void swap(int* a, int* b) {
  int const tmp = *a;
  *a = *b;
  *b = tmp;
}

int main(void) {
  int i = 3, j = 2;
  printf("%d,%d\n", i, j); // 3,2
  swap(&i, &j);
  printf("%d,%d\n", i, j); // 2,3
  return 0;
}
\end{lstlisting}

Rappel : c’est exactement pour ça qu’on écrit \texttt{scanf("\%d", \&x);} 

\subsection{Optimisation : passer par adresse sans vouloir modifier}
Passer un pointeur peut éviter une copie (structures/objets volumineux), mais si ce
n’est \emph{pas} un vrai « passage par référence », on protège avec \texttt{const}. 

\begin{lstlisting}[language=C]
typedef struct { double a[3][3]; } Matrice;

Matrice addition(Matrice const* m1, Matrice const* m2);
\end{lstlisting}

\subsection{const : pointeur sur const vs pointeur const}
Règle : \texttt{const} s’applique au type juste avant (ou juste après s’il est au début). 

\begin{itemize}
  \item \texttt{T const* p} (ou \texttt{const T* p}) : \textbf{pointeur vers const}
  \begin{itemize}
    \item on peut changer \texttt{p} (le faire pointer ailleurs),
    \item on ne peut pas modifier \texttt{*p}.
  \end{itemize}
  \item \texttt{T* const p = \&x} : \textbf{pointeur constant}
  \begin{itemize}
    \item on ne peut pas changer \texttt{p},
    \item on peut modifier \texttt{*p}.
  \end{itemize}
\end{itemize}

\begin{lstlisting}[language=C]
int i = 2, j = 3;

int const* p1 = &i;  // pointeur sur int const
int* const p2 = &i;  // pointeur const sur int

// *p1 = 5; // interdit
*p2 = 5;     // OK

p1 = &j;     // OK
// p2 = &j;  // interdit
\end{lstlisting}

\subsection{Pointeurs vs références (conceptuellement)}
Le C n’a pas de références (comme en C++), mais le cours compare : 
\begin{itemize}
  \item une \textbf{référence} est un alias (un autre nom), toujours initialisée, jamais nulle,
        ne peut pas changer d’objet référencé.
  \item « Pour faire simple » : une référence $\approx$ un pointeur \texttt{* const} toujours valide.
  \item Adage (dans les langages qui ont les deux) : utiliser références si possible, pointeurs si nécessaire.
\end{itemize}

\subsection{Allocation mémoire : pile vs tas}
Deux façons d’allouer : 
\begin{enumerate}
  \item \textbf{Allocation statique} (variables) : taille connue à la compilation, sur la \textbf{pile} (stack) pour les variables locales.
  \item \textbf{Allocation dynamique} : pendant l’exécution, sur le \textbf{tas} (heap).
\end{enumerate}

Mémoire virtuelle typique d’un processus : texte (code+constantes), données (globales/statiques), tas, pile.

\subsection{malloc, calloc, sizeof, size\_t}
\subsubsection{malloc}
\texttt{p = malloc(taille);} réserve \texttt{taille} octets et retourne l’adresse, ou \texttt{NULL} si échec. 

\begin{lstlisting}[language=C]
#include <stdlib.h>

double* px = malloc(sizeof(double)); // place pour un double
\end{lstlisting}

\texttt{sizeof(expr)} retourne un \texttt{size\_t}; format printf : \texttt{\%zu}.

\subsubsection{calloc (préféré pour allocations multiples)}
\texttt{calloc(nb, taille\_el)} :
\begin{itemize}
  \item alloue \textbf{nb} éléments consécutifs,
  \item \textbf{initialise à 0},
  \item et protège contre certains \textbf{overflows} de multiplication. 
\end{itemize}

\begin{lstlisting}[language=C]
double* a = calloc(3, sizeof(double)); // 3 doubles, init a 0
\end{lstlisting}

Le cours insiste : \texttt{malloc(n * sizeof(T))} peut overflow et sous-allouer $\rightarrow$ puis \textbf{buffer overflow} en accédant à \texttt{p[i]}.

\subsubsection{Initialiser la mémoire : memset}
Avec malloc, rien n’est initialisé. Pour initialiser :
\texttt{memset(ptr, valeur, taille)} (dans \texttt{string.h}).

\begin{lstlisting}[language=C]
#include <string.h>

struct Machin { int x; double y; };
struct Machin m;
memset(&m, 0, sizeof(m));
\end{lstlisting}

\subsubsection{Toujours tester l’allocation}
\begin{lstlisting}[language=C]
p = calloc(n, sizeof(*p));
if (p == NULL) {
  /* gestion erreur + return */
}
\end{lstlisting}

\subsection{free : libérer la mémoire}
\texttt{free(p)} libère la zone allouée; ensuite \textbf{ne plus accéder} à cette zone.   
Conseil de prudence : mettre \texttt{p = NULL;} après un free.

\begin{lstlisting}[language=C]
free(p);
p = NULL;
\end{lstlisting}

\paragraph{Règle absolue.}
Toute zone allouée par \texttt{malloc/calloc} doit être libérée par un \texttt{free} correspondant. 

\subsection{Toujours allouer avant d’utiliser}
Déréférencer un pointeur non initialisé/non alloué cause souvent \textit{segfault} / \textit{bus error}. 

\begin{lstlisting}[language=C]
int* px;     // non initialise
*px = 20;    // ERREUR : segfault probable
\end{lstlisting}

Conseil : toujours initialiser, typiquement à \texttt{NULL} tant qu’on n’a pas alloué. 

\subsection{Segmentation fault vs Bus error}
Les deux = accès mémoire interdit, mais \textit{bus error} peut indiquer que le matériel
a détecté un problème que le noyau n’a pas détecté « proprement ».

\subsection{Bonnes pratiques (récapitulatif du cours)}
Règles explicitement listées : 
\begin{itemize}
  \item Toute zone \texttt{malloc/calloc} doit avoir un \texttt{free} correspondant (et celui qui alloue doit libérer).
  \item Pour allocations multiples : préférer \texttt{calloc} à \texttt{malloc}.
  \item Tester systématiquement \texttt{malloc/calloc}.
  \item Toujours initialiser les pointeurs (NULL si inconnu).
  \item Toujours initialiser la mémoire utilisée (ex. \texttt{memset}).
  \item Mettre \texttt{ptr = NULL;} après chaque \texttt{free}.
  \item Mettre \texttt{const} dans les « faux passages par référence » (optimisation).
  \item Utiliser outils : options compilateur, debugger, valgrind, AddressSanitizer, etc.
\end{itemize}

\subsection{Tableaux dynamiques : pourquoi et interface}
Un tableau dynamique = collection homogène à accès direct dont la taille n’est pas fixée a priori. 

Interface typique :
\begin{itemize}
  \item accéder/modifier un élément
  \item push/pop en fin
  \item tester vide
  \item itérer
\end{itemize}
En C, pas de \texttt{ArrayList}/\texttt{vector} standard $\rightarrow$ à implémenter soi-même.

\subsection{Exemple : implémenter un \texttt{vector} en C}
\subsubsection{Structure de données}
Idée : stocker
\begin{itemize}
  \item \texttt{size} : nb d’éléments utilisés,
  \item \texttt{allocated} : capacité allouée,
  \item \texttt{content} : tableau dynamique (heap). 
\end{itemize}

\begin{lstlisting}[language=C]
typedef int type_el;

typedef struct {
  size_t size;       // nb d'elements utilises
  size_t allocated;  // capacite allouee
  type_el* content;  // zone allouee dynamiquement
} vector;
\end{lstlisting}

\subsubsection{realloc : redimensionner}
Propriétés clés : 
\begin{itemize}
  \item augmente/diminue la taille
  \item peut déplacer la zone (copie), conserve les données
  \item si échec : retourne \texttt{NULL} et l’ancienne zone reste inchangée
  \item la mémoire supplémentaire n’est pas initialisée
  \item \texttt{realloc(NULL, n)} $\equiv$ \texttt{malloc(n)}
  \item \texttt{realloc(p, 0)} $\equiv$ \texttt{free(p)}
\end{itemize}

\subsubsection{Construction / initialisation}
Le cours propose une construction « atomique » : on construit un \texttt{result} local, on alloue, puis on fait \texttt{*v = result;} à la fin.
\begin{lstlisting}[language=C]
#include <stdlib.h>

#define VECTOR_PADDING 128

vector* vector_construct(vector* v) {
  if (v != NULL) {
    vector result = (vector){ 0, 0, NULL };
    result.content = calloc(VECTOR_PADDING, sizeof(type_el));
    if (result.content == NULL) return NULL;
    result.allocated = VECTOR_PADDING;
    *v = result; // ecriture "atomique"
  }
  return v;
}
\end{lstlisting}

\subsubsection{Destruction}
Comme on fournit une fonction qui alloue, on doit fournir une fonction qui libère. 

\begin{lstlisting}[language=C]
#include <stdlib.h>

void vector_delete(vector* v) {
  if (v != NULL && v->content != NULL) {
    free(v->content);
    v->content = NULL;
    v->size = 0;
    v->allocated = 0;
  }
}
\end{lstlisting}

Rappel : \texttt{x->y} $\equiv$ \texttt{(*x).y}. 

\subsubsection{Ajouter un élément (push)}
On agrandit si nécessaire, puis on écrit et on incrémente. 

\begin{lstlisting}[language=C]
size_t vector_push(vector* v, type_el val) {
  if (v != NULL) {
    while (v->size >= v->allocated) {
      if (vector_enlarge(v) == NULL) return 0;
    }
    v->content[v->size] = val;
    ++(v->size);
    return v->size;
  }
  return 0;
}
\end{lstlisting}

\subsubsection{Agrandir (enlarge) et éviter overflow}
Le cours montre un test basé sur \texttt{SIZE\_MAX} pour éviter overflow avant de calculer \texttt{allocated * sizeof(type\_el)}.

\begin{lstlisting}[language=C]
#include <stdint.h>
#include <stdlib.h>

#ifndef SIZE_MAX
#define SIZE_MAX (~(size_t)0)
#endif

vector* vector_enlarge(vector* v) {
  if (v != NULL) {
    vector result = *v;
    result.allocated += VECTOR_PADDING;

    if (result.allocated > SIZE_MAX / sizeof(type_el)) return NULL;

    result.content = realloc(result.content, result.allocated * sizeof(type_el));
    if (result.content == NULL) return NULL; // v inchangé

    *v = result;
  }
  return v;
}
\end{lstlisting}

\subsection{Récapitulatif (formules à retenir)}
\begin{itemize}
  \item Déclaration : \texttt{type* p;}
  \item Adresse : \texttt{\&x}
  \item Contenu pointé : \texttt{*p}
  \item Pointeur sur const : \texttt{type const* p;}
  \item Pointeur const : \texttt{type* const p = ...;}
  \item Allocation : \texttt{malloc}, \texttt{calloc}, \texttt{realloc}
  \item Libération : \texttt{free(p);}
  \item Pointeur de fonction : \texttt{ret (*pf)(args...)}
\end{itemize}

